% --- SLIDE A (GỘP TỪ SLIDE 4 & 5) ---
\begin{frame}
    \frametitle{Giới thiệu tổng quan}
    \SubTitle{Bối cảnh và Vấn đề}
    \begin{BulletTight}
        \item \textbf{Bối cảnh:}
        \begin{BulletTight}
            \item Kiểm thử là giai đoạn then chốt đảm bảo chất lượng phần mềm.
            \item Chi phí sửa lỗi sau phát hành cao gấp 10–100 lần phát hiện sớm.\src{1}{Boehm, B. W. (1981). \textit{Software engineering economics}.}
            \item Phần mềm phức tạp, thay đổi nhanh \textrightarrow{} Tự động hóa là tất yếu.
        \end{BulletTight}
        \item \textbf{Vấn đề (Trước LLM):}
        \begin{BulletTight}
            \item Các phương pháp truyền thống (kịch bản, học tăng cường [2]) còn hạn chế: \src{2}{Nguyen, A. et al. (2019). Prioritizing automated UI tests using RL.}
            \item Khó thích ứng khi phần mềm thay đổi.
            \item Thiếu khả năng hiểu ngữ nghĩa và bối cảnh kiểm thử.
            \item Độ bao phủ logic thấp, dễ bỏ sót lỗi quan trọng.
        \end{BulletTight}
    \end{BulletTight}
\end{frame}

% --- SLIDE B (GỘP TỪ SLIDE 6 & 7) ---
\begin{frame}
    \frametitle{Giới thiệu tổng quan}
    \SubTitle{Tiềm năng & Hạn chế của LLM (Đơn tác tử)}
    \begin{BulletTight}
        \item \textbf{Tiềm năng (Đơn tác tử với LLM):}
        \begin{BulletTight}
            \item LLM tạo bước tiến lớn trong tự động hóa kiểm thử.
            \item Hiểu ngữ cảnh, sinh test ngữ nghĩa, tự sửa lỗi (InferFix).\src{3}{Jin, M. et al. (2023). InferFix: End-to-end program repair with LLMs.}
            \item Mở rộng kiểm thử UI theo hành vi người dùng (Guardian).\src{4}{Ran, D. et al. (2024). Guardian: LLM-based UI exploration.}
        \end{BulletTight}
        \item \textbf{Hạn chế (Đơn tác tử):}
        \begin{BulletTight}
            \item Một tác tử xử lý toàn bộ quy trình \textrightarrow{} Thiếu phối hợp.\src{5}{Fan, Y. et al. (2024). Can cooperative multi-agent RL boost automatic web testing?}
            \item Dễ lặp lại hành động, bỏ sót lỗi.
            \item Dễ "ảo giác" (Hallucination) mục tiêu.\src{6}{Xu, Q. et al. (2025). Hallucination to consensus: Multi-agent LLMs.}
        \end{BulletTight}
        \item[\raise0.2ex\hbox{\color{KHTN_BLUE!85}\scriptsize\faArrowRight}] Chưa phù hợp cho kiểm thử quy mô lớn, đa nhiệm.
    \end{BulletTight}
\end{frame}

% --- SLIDE C (GỘP TỪ SLIDE 8, 9 & 10) ---
\begin{frame}
    \frametitle{Giới thiệu tổng quan}
    \SubTitle{Hướng tiếp cận Đa tác tử & Thách thức}
    \begin{BulletTight}
        \item \textbf{Giải pháp (Hệ Đa tác tử):}
        \begin{BulletTight}
            \item Mô phỏng quy trình phát triển đa vai trò (MetaGPT).\src{7}{Hong, S. et al. (2024). Meta programming for multi-agent collaboration.}
            \item Phối hợp nhiều tác tử tăng bao phủ và hiệu quả.\src{8}{Feng, S. et al. (2025). Can cooperative multi-agent RL boost automatic web testing?}
            \item Cộng tác & đồng thuận giúp giảm trùng lặp, giảm "ảo giác".\src{9}{He, J. et al. (2025). LLM-based multi-agent systems.}\src{10}{Xu, Q. et al. (2025). Hallucination to consensus...}
        \end{BulletTight}
        \item \textbf{Thách thức chính (Cần giải quyết):}
        \begin{BulletTight}
            \item \textbf{Ảo giác (Hallucination):} Tác tử sinh phản hồi sai lệch [5].
            \item \textbf{Giao tiếp / Phối hợp:} Dễ mất ngữ cảnh, lan truyền lỗi [9].
            \item \textbf{HITL (Con người):} Giám sát tốn chi phí, phản hồi thiên kiến [10].
            \item \textbf{Hiệu năng:} Tốn tài nguyên, tăng độ trễ khi chạy nhiều LLM [7].
        \end{BulletTight}
    \end{BulletTight}
\end{frame}